
\documentclass[11pt]{article}

\usepackage[margin=1in]{geometry}
\usepackage{graphicx}
\usepackage{float}
\usepackage{booktabs}
\usepackage{longtable}
\usepackage{array}
\usepackage{caption}
\usepackage{hyperref}
\usepackage{amsmath}

\title{Combined Results and Discussion for Encoder Models}
\author{}
\date{}

\begin{document}

\maketitle

\section{Gemma}

\subsection{Exploratory Analysis}

Figure~\ref{fig:bert_line} illustrates the relationship between chunk size and BERTScore F1 across different maximum token settings. A clear downward trend is observed, where smaller chunk sizes (e.g., 128) consistently achieve higher performance, while larger chunk sizes (e.g., 1024) show degraded performance. This suggests that smaller chunks preserve semantic coherence more effectively during question-answer generation.

\begin{figure}[h]
    \centering
    \includegraphics[width=0.8\linewidth]{lineplot_bert_f1.png}
    \caption{BERTScore F1 by Chunk Size and Max Tokens}
    \label{fig:bert_line}
\end{figure}

The boxplot in Figure~\ref{fig:bert_box} further confirms this pattern, showing tighter distributions and higher median performance for smaller chunk sizes. Larger chunk sizes exhibit lower central tendency and increased variability, indicating less stable model behavior.

\begin{figure}[h]
    \centering
    \includegraphics[width=0.8\linewidth]{boxplot_bert_f1.png}
    \caption{Distribution of BERTScore F1 Across Chunk Sizes}
    \label{fig:bert_box}
\end{figure}

\subsection{Correlation Analysis and Multicollinearity}

The correlation matrix in Figure~\ref{fig:corr_heatmap} reveals extremely high correlations among ROUGE metrics (approaching 1.00) and strong correlations with BERTScore and STS. This indicates substantial multicollinearity among dependent variables.

\begin{figure}[h]
    \centering
    \includegraphics[width=0.8\linewidth]{correlation_heatmap.png}
    \caption{Correlation Heatmap of Evaluation Metrics}
    \label{fig:corr_heatmap}
\end{figure}

Such high correlations suggest that many of the evaluation metrics capture similar underlying information, which can violate the independence assumptions required for parametric multivariate analysis.

\subsection{Assessment of MANOVA Assumptions}

Several diagnostic tests were conducted to evaluate the suitability of MANOVA. The Henze-Zirkler test indicated a significant violation of multivariate normality ($p < 0.001$). Additionally, Box’s M test showed that the assumption of homogeneity of covariance matrices was not satisfied ($p < 0.001$). Outlier analysis further revealed that approximately 25\% of observations were classified as multivariate outliers.

The Q-Q plot of Mahalanobis distances (Figure~\ref{fig:qq_plot}) demonstrates substantial deviation from the theoretical distribution, further confirming the violation of normality assumptions.

\begin{figure}[h]
    \centering
    \includegraphics[width=0.8\linewidth]{qq_plot_mahalanobis.png}
    \caption{Q-Q Plot of Mahalanobis Distances}
    \label{fig:qq_plot}
\end{figure}

Due to these violations, the use of MANOVA is not considered reliable for inference in this context.

\subsection{PERMANOVA Results}

To address the limitations of MANOVA, PERMANOVA was employed as a non-parametric alternative. The results indicate a statistically significant effect of chunk size on the combined evaluation metrics (Pseudo-F = 35.15, $p = 0.001$). 

This finding confirms that chunk size has a meaningful impact on model performance across all evaluated metrics. Unlike MANOVA, PERMANOVA does not rely on strict distributional assumptions, making it more appropriate for this dataset.

\subsection{Dimensionality Reduction Analysis}

Principal Component Analysis (PCA) was used to visualize group separation in reduced dimensional space. As shown in Figure~\ref{fig:pca}, distinct clustering patterns emerge for different chunk sizes, indicating that the groups are well-separated based on their combined metric profiles.

\begin{figure}[h]
    \centering
    \includegraphics[width=0.8\linewidth]{pca_scatter.png}
    \caption{PCA of Dependent Variables by Chunk Size}
    \label{fig:pca}
\end{figure}

The separation observed in PCA space aligns with the PERMANOVA results, reinforcing the conclusion that chunk size significantly influences model outputs.

\subsection{Discussion}

The results demonstrate that chunk size is a critical factor in determining the quality of generated question-answer pairs. Smaller chunk sizes consistently yield higher performance across all evaluation metrics, likely due to improved contextual coherence and reduced information overload during generation.

The strong multicollinearity among evaluation metrics suggests that they measure overlapping aspects of performance, which may reduce the need for multiple redundant metrics in future analyses.

The violation of MANOVA assumptions highlights the importance of validating statistical assumptions before applying parametric methods. In this study, PERMANOVA provided a robust alternative, enabling reliable inference despite the presence of non-normality, outliers, and unequal covariance structures.

Overall, the findings suggest that optimizing chunk size is essential for improving the effectiveness of retrieval-augmented generation systems, and that non-parametric multivariate methods should be preferred when dealing with complex, high-dimensional evaluation data.




\section{Mistrial}

\subsection{Exploratory Analysis}

Figure~\ref{fig:bert_line_mistral} shows the relationship between chunk size and BERTScore F1 for the Mistral model. A consistent downward trend is observed, where smaller chunk sizes (128) achieve the highest performance, while larger chunk sizes (1024) result in reduced performance. This indicates that smaller chunks preserve contextual coherence more effectively during generation.

\begin{figure}[h]
    \centering
    \includegraphics[width=0.8\linewidth]{lineplot_bert_f1.png}
    \caption{BERTScore F1 by Chunk Size and Max Tokens (Mistral)}
    \label{fig:bert_line_mistral}
\end{figure}

The distribution of performance across chunk sizes (Figure~\ref{fig:bert_box_mistral}) further confirms this pattern, with smaller chunk sizes exhibiting higher and more stable scores, while larger chunk sizes show lower central tendency.

\begin{figure}[h]
    \centering
    \includegraphics[width=0.8\linewidth]{boxplot_bert_f1.png}
    \caption{Distribution of BERTScore F1 Across Chunk Sizes (Mistral)}
    \label{fig:bert_box_mistral}
\end{figure}

\subsection{Correlation Analysis and Multicollinearity}

The correlation matrix (Figure~\ref{fig:corr_heatmap_mistral}) reveals extremely high correlations among ROUGE metrics (approaching 1.00) and strong correlations with BERTScore and STS metrics. This indicates severe multicollinearity among dependent variables.

\begin{figure}[h]
    \centering
    \includegraphics[width=0.8\linewidth]{correlation_heatmap.png}
    \caption{Correlation Heatmap of Evaluation Metrics (Mistral)}
    \label{fig:corr_heatmap_mistral}
\end{figure}

Such high correlations suggest redundancy among evaluation metrics, which violates the independence assumptions required for parametric multivariate analysis.

\subsection{Assessment of MANOVA Assumptions}

Statistical diagnostics revealed multiple violations of MANOVA assumptions. The Henze-Zirkler test indicated a significant deviation from multivariate normality ($p < 0.001$) :contentReference[oaicite:0]{index=0}. Additionally, Box’s M test confirmed that covariance matrices were not equal across groups ($p = 0.0$) :contentReference[oaicite:1]{index=1}. 

Outlier analysis showed that approximately 25\% of observations were classified as multivariate outliers :contentReference[oaicite:2]{index=2}. The Q-Q plot of Mahalanobis distances (Figure~\ref{fig:qq_plot_mistral}) further demonstrates strong deviation from the theoretical distribution, confirming the violation of normality assumptions.

\begin{figure}[h]
    \centering
    \includegraphics[width=0.8\linewidth]{qq_plot_mahalanobis.png}
    \caption{Q-Q Plot of Mahalanobis Distances (Mistral)}
    \label{fig:qq_plot_mistral}
\end{figure}

Due to these violations, the use of MANOVA is not considered reliable for inference.

\subsection{PERMANOVA Results}

To address these violations, PERMANOVA was used as a non-parametric alternative. The results indicate a highly significant effect of chunk size on the combined evaluation metrics (Pseudo-F = 124.07, $p = 0.001$) :contentReference[oaicite:3]{index=3}.

The magnitude of the pseudo-F statistic suggests strong separation between chunk size groups, indicating that chunk size has a substantial impact on model performance across all evaluation metrics.

\subsection{Dimensionality Reduction Analysis}

Principal Component Analysis (PCA) was used to visualize group separation. As shown in Figure~\ref{fig:pca_mistral}, distinct clustering patterns are observed across chunk sizes, indicating that groups are well-separated in reduced dimensional space.

\begin{figure}[h]
    \centering
    \includegraphics[width=0.8\linewidth]{pca_scatter.png}
    \caption{PCA of Dependent Variables by Chunk Size (Mistral)}
    \label{fig:pca_mistral}
\end{figure}

The separation observed in PCA space is consistent with PERMANOVA results, further confirming the effect of chunk size on model outputs.

\subsection{Discussion}

The results demonstrate that chunk size significantly influences the performance of the Mistral model. Smaller chunk sizes consistently produce higher evaluation scores, suggesting that limiting context length improves semantic coherence and generation quality.

Compared to the previous model, the PERMANOVA pseudo-F value is substantially higher (124.07), indicating stronger group separation and a more pronounced effect of chunk size on performance. This suggests that the Mistral model is more sensitive to chunk size variations.

The presence of strong multicollinearity among evaluation metrics indicates that these metrics capture overlapping aspects of performance. Future work may consider dimensionality reduction or metric selection to reduce redundancy.

The violation of MANOVA assumptions highlights the importance of validating statistical conditions before applying parametric tests. In this study, PERMANOVA provided a robust alternative, enabling reliable inference despite non-normality, outliers, and unequal covariance structures.

Overall, these findings reinforce the importance of chunk size optimization in retrieval-augmented generation systems and demonstrate that non-parametric multivariate methods are more appropriate for analyzing high-dimensional evaluation data in this context.





\section{Phi Model}

\subsection{Exploratory Analysis}

Figure~\ref{fig:bert_line_phi} shows the relationship between chunk size and BERTScore F1 for the Phi model. Similar to previous models, performance decreases as chunk size increases. However, the decline is less steep compared to the Mistral model, indicating a weaker sensitivity to chunk size.

\begin{figure}[h]
    \centering
    \includegraphics[width=0.8\linewidth]{lineplot_bert_f1.png}
    \caption{BERTScore F1 by Chunk Size and Max Tokens (Phi)}
    \label{fig:bert_line_phi}
\end{figure}

The boxplot (Figure~\ref{fig:bert_box_phi}) confirms that smaller chunk sizes achieve higher performance, though the spread across chunk sizes is more uniform compared to previous models.

\begin{figure}[h]
    \centering
    \includegraphics[width=0.8\linewidth]{boxplot_bert_f1.png}
    \caption{Distribution of BERTScore F1 Across Chunk Sizes (Phi)}
    \label{fig:bert_box_phi}
\end{figure}

\subsection{Correlation Analysis and Multicollinearity}

The correlation heatmap (Figure~\ref{fig:corr_heatmap_phi}) shows strong correlations among ROUGE metrics and moderate correlations with BERTScore precision. Compared to previous models, BERTScore precision exhibits weaker correlation with other metrics, suggesting slightly reduced redundancy.

\begin{figure}[h]
    \centering
    \includegraphics[width=0.8\linewidth]{correlation_heatmap.png}
    \caption{Correlation Heatmap of Evaluation Metrics (Phi)}
    \label{fig:corr_heatmap_phi}
\end{figure}

Despite this, multicollinearity remains present and continues to violate MANOVA assumptions.

\subsection{Assessment of MANOVA Assumptions}

Diagnostic tests indicate that MANOVA assumptions are violated. The Henze-Zirkler test confirms significant deviation from multivariate normality ($p < 0.001$) :contentReference[oaicite:0]{index=0}. Box’s M test indicates unequal covariance matrices across groups ($p = 0.0$) :contentReference[oaicite:1]{index=1}.

Additionally, 25\% of observations were identified as multivariate outliers :contentReference[oaicite:2]{index=2}. The Q-Q plot of Mahalanobis distances (Figure~\ref{fig:qq_plot_phi}) further confirms deviation from theoretical normality.

\begin{figure}[h]
    \centering
    \includegraphics[width=0.8\linewidth]{qq_plot_mahalanobis.png}
    \caption{Q-Q Plot of Mahalanobis Distances (Phi)}
    \label{fig:qq_plot_phi}
\end{figure}

These violations render MANOVA unreliable for inference.

\subsection{PERMANOVA Results}

PERMANOVA results indicate a statistically significant effect of chunk size on the combined evaluation metrics (Pseudo-F = 13.63, $p = 0.001$) :contentReference[oaicite:3]{index=3}.

While the effect is statistically significant, the pseudo-F value is substantially lower than that observed in the Mistral model, suggesting weaker separation between chunk size groups.

\subsection{Dimensionality Reduction Analysis}

The PCA plot (Figure~\ref{fig:pca_phi}) shows partial clustering of chunk sizes, but with greater overlap between groups compared to previous models. This indicates that chunk size still influences performance, but the separation is less distinct.

\begin{figure}[h]
    \centering
    \includegraphics[width=0.8\linewidth]{pca_scatter.png}
    \caption{PCA of Dependent Variables by Chunk Size (Phi)}
    \label{fig:pca_phi}
\end{figure}

\subsection{Discussion}

The results indicate that chunk size has a statistically significant effect on the performance of the Phi model. However, the magnitude of this effect is smaller compared to other models, as reflected by the lower PERMANOVA pseudo-F value.

This suggests that the Phi model is less sensitive to chunk size variations, exhibiting more stable performance across different chunk sizes. While smaller chunk sizes still yield the best results, the performance degradation at larger chunk sizes is less severe than observed in the Mistral model.

The presence of multicollinearity and violations of MANOVA assumptions reinforces the need for non-parametric approaches such as PERMANOVA. The reduced separation observed in PCA space further supports the conclusion that chunk size has a weaker influence on this model.

Overall, these findings suggest that while chunk size optimization remains important, the Phi model demonstrates greater robustness to chunk size changes compared to other evaluated models.

\section{Qwen Model}

\subsection{Exploratory Analysis}

Figure~\ref{fig:bert_line_qwen} shows the relationship between chunk size and BERTScore F1 for the Qwen model. A strong and consistent decline in performance is observed as chunk size increases. Smaller chunk sizes (128) achieve the highest scores, while larger chunk sizes (1024) result in noticeably lower performance.

\begin{figure}[h]
    \centering
    \includegraphics[width=0.8\linewidth]{lineplot_bert_f1.png}
    \caption{BERTScore F1 by Chunk Size and Max Tokens (Qwen)}
    \label{fig:bert_line_qwen}
\end{figure}

The boxplot (Figure~\ref{fig:bert_box_qwen}) confirms this trend, showing higher median values and tighter distributions for smaller chunk sizes, while larger chunk sizes exhibit reduced performance.

\begin{figure}[h]
    \centering
    \includegraphics[width=0.8\linewidth]{boxplot_bert_f1.png}
    \caption{Distribution of BERTScore F1 Across Chunk Sizes (Qwen)}
    \label{fig:bert_box_qwen}
\end{figure}

\subsection{Correlation Analysis and Multicollinearity}

The correlation heatmap (Figure~\ref{fig:corr_heatmap_qwen}) shows extremely high correlations among ROUGE metrics and strong correlations with BERTScore and STS. This indicates substantial multicollinearity among dependent variables.

\begin{figure}[h]
    \centering
    \includegraphics[width=0.8\linewidth]{correlation_heatmap.png}
    \caption{Correlation Heatmap of Evaluation Metrics (Qwen)}
    \label{fig:corr_heatmap_qwen}
\end{figure}

\subsection{Assessment of MANOVA Assumptions}

Unlike previous models, the Henze-Zirkler test indicates that multivariate normality is not violated ($p = 0.091$) :contentReference[oaicite:0]{index=0}. However, other assumptions remain violated. Box’s M test confirms unequal covariance matrices across groups ($p = 0.0$) :contentReference[oaicite:1]{index=1}, and 25\% of observations are identified as multivariate outliers :contentReference[oaicite:2]{index=2}.

The Q-Q plot of Mahalanobis distances (Figure~\ref{fig:qq_plot_qwen}) shows deviation from the theoretical distribution, suggesting that the presence of outliers still affects the distributional assumptions.

\begin{figure}[h]
    \centering
    \includegraphics[width=0.8\linewidth]{qq_plot_mahalanobis.png}
    \caption{Q-Q Plot of Mahalanobis Distances (Qwen)}
    \label{fig:qq_plot_qwen}
\end{figure}

Although normality is satisfied, the violation of covariance homogeneity and presence of outliers limit the reliability of MANOVA.

\subsection{PERMANOVA Results}

PERMANOVA results indicate a highly significant effect of chunk size on the combined evaluation metrics (Pseudo-F = 220.85, $p = 0.001$) :contentReference[oaicite:3]{index=3}.

This is the largest pseudo-F value observed across all evaluated models, indicating extremely strong separation between chunk size groups.

\subsection{Dimensionality Reduction Analysis}

The PCA plot (Figure~\ref{fig:pca_qwen}) shows very clear separation between chunk size groups, with minimal overlap. This indicates that chunk size has a strong and consistent impact on model outputs.

\begin{figure}[h]
    \centering
    \includegraphics[width=0.8\linewidth]{pca_scatter.png}
    \caption{PCA of Dependent Variables by Chunk Size (Qwen)}
    \label{fig:pca_qwen}
\end{figure}

\subsection{Discussion}

The results demonstrate that chunk size has a substantial impact on the performance of the Qwen model. The very high PERMANOVA pseudo-F value (220.85) indicates that chunk size strongly influences the distribution of evaluation metrics.

Compared to other models, Qwen exhibits the highest sensitivity to chunk size, suggesting that careful tuning of chunk size is critical for optimal performance. Smaller chunk sizes consistently yield the best results, while performance declines as chunk size increases.

Although multivariate normality is satisfied for this model, other MANOVA assumptions remain violated, reinforcing the use of PERMANOVA as a more robust alternative.

The clear separation observed in PCA space further confirms that chunk size plays a dominant role in determining model performance. These findings suggest that Qwen benefits significantly from smaller chunk sizes and that optimization of chunking strategy is essential when deploying this model in retrieval-augmented generation systems.



\subsection{Cross-Model Comparison}

A clear ranking emerges when comparing the PERMANOVA statistics across models:

\begin{table}[H]
\centering
\caption{Ranking of models by sensitivity to chunk size}
\label{tab:model_ranking}
\begin{tabular}{lc}
\toprule
Model & PERMANOVA pseudo-$F$ \\
\midrule
Qwen    & 220.85 \\
Mistral & 124.07 \\
Gemma   & 35.15 \\
Phi     & 13.63 \\
\bottomrule
\end{tabular}
\end{table}

These results indicate that different model architectures exhibit substantially different sensitivity to chunk size. Qwen and Mistral are highly sensitive, meaning that retrieval chunking strategy has a strong influence on their downstream performance. In contrast, Phi appears to be comparatively robust to chunk-size variation, while Gemma occupies an intermediate position.

\subsection{Overall Discussion}

Across all four models, smaller chunk sizes consistently yielded better performance than larger chunk sizes. This pattern likely reflects improved semantic coherence and reduced information dilution when smaller chunks are provided as input to retrieval-augmented generation systems.

However, the degree of chunk-size sensitivity differs substantially by model. This suggests that chunk-size optimization should not be treated as a one-size-fits-all design choice. Instead, chunking strategy should be tuned jointly with model selection.

Finally, the repeated violation of MANOVA assumptions across models, particularly covariance heterogeneity, outlier prevalence, and multicollinearity, demonstrates the importance of validating statistical assumptions before relying on parametric multivariate methods. In this study, PERMANOVA provided a robust and consistent alternative for detecting group-level differences in multivariate evaluation outcomes.

\end{document}
